\documentclass[11pt]{article}

\usepackage{series}   % provided style file
\usepackage{amsmath,amssymb,amsthm}
\usepackage{mathtools}
\usepackage{hyperref}

\title{Deterministic Projection, Diffusive Limits, and\\
Knee--Like Threshold Transitions}

\author{Peter Nero}

\date{January, 2026}

\begin{document}
\maketitle

\begin{abstract}
Abrupt threshold phenomena---such as collapse--like transitions, sharp metastable escape, or rapid post--selection relaxation---are often modeled by introducing intrinsic stochasticity, explicit discontinuities, or ad hoc collapse rules. In this paper we show that none of these are necessary. We construct a fully deterministic, finite--dimensional dynamical system whose \emph{non--injective projection} onto an observed variable produces stochastic effective behavior and a sharp, finite--strength ``knee'' transition in first--passage statistics. 

Using rigorous deterministic homogenization for chaotic fast subsystems, we derive an Ornstein--Uhlenbeck limit for the projected coordinate and prove convergence of exit probabilities to those of the limiting diffusion. The knee arises when an effective stability margin crosses zero, separating a noise--activated escape regime from deterministic expulsion. We further show that such knee behavior cannot be reproduced by fixed, protocol--independent linear Markov generators without additional state--dependent structure. 

The result provides a minimal, explicit stress test for projection--based descriptions and clarifies how sudden effective transitions can arise from smooth, invertible dynamics.
\end{abstract}

\section{Introduction}

\subsection{Motivation}

Many physical phenomena exhibit abrupt effective transitions: quantum measurement outcomes appear suddenly discrete; metastable systems escape confinement sharply as parameters vary; horizons and singularities mark apparent breakdowns of classical spacetime description; and inflationary phases display rapid relaxation following critical events. Standard explanations typically invoke one or more of the following mechanisms:
\begin{itemize}
\item intrinsic stochasticity or noise,
\item explicit discontinuities or collapse postulates,
\item thermodynamic limits with infinite degrees of freedom,
\item or special dynamical ingredients added by hand.
\end{itemize}
While these approaches are often effective phenomenologically, they obscure a more basic structural question:

\medskip
\emph{Can sharp effective transitions arise generically from smooth, deterministic, invertible dynamics once one admits non--injective projection onto observed variables?}
\medskip

This question is not merely interpretive. In many settings the observer has access only to a reduced description, while the full dynamics---including unobserved degrees of freedom---remains deterministic and well defined. The reduced description is therefore obtained by a \emph{projection} that identifies many microscopic states as equivalent. Such projections are necessarily non--injective, and it is natural to ask whether they can force qualitative changes in effective behavior without modifying the underlying dynamics.

\subsection{Main idea and result}

The central result of this paper is that the answer is \emph{yes}. We exhibit a minimal deterministic system in which:
\begin{enumerate}
\item the full ``upstairs'' dynamics is smooth, deterministic, and invertible;
\item the observed ``downstairs'' variable is obtained by a non--injective projection;
\item deterministic homogenization yields a diffusive effective limit for the projected variable;
\item a finite control parameter produces a sharp \emph{knee--like crossover} in first--passage statistics.
\end{enumerate}

The knee occurs when an effective stability coefficient crosses zero. On one side of the critical value, boundary crossing is noise--activated and strongly suppressed; on the other, deterministic drift dominates and escape occurs rapidly. The transition is continuous but sharply bent, with a width controlled by the effective diffusion scale. No thermodynamic limit, intrinsic randomness, or discontinuous dynamics is required.

\subsection{Why this matters}

Threshold behavior of this kind is often discussed in quantum measurement, collapse models, and metastability theory, but usually by introducing additional structure at the effective level. By contrast, the present construction shows that:
\begin{itemize}
\item stochastic effective behavior can emerge from deterministic chaos via projection alone;
\item sharp threshold transitions can arise from loss of stability in the effective encoding;
\item the location and width of the knee are determined by stability margins, not by ad hoc rules.
\end{itemize}

This makes the knee phenomenon a natural stress test for any framework that treats projection as fundamental rather than incidental. In particular, it allows one to distinguish between explanations that rely on fixed linear generators and those that incorporate state-- or protocol--dependent effective dynamics induced by discarded degrees of freedom.

\subsection{Structure of the paper}

Section~2 introduces the explicit deterministic toy model and the projection defining the observed variable. In Section~3 we derive a rigorous diffusive limit using deterministic homogenization and obtain an Ornstein--Uhlenbeck process near the basin boundary. Section~4 establishes the existence of a finite critical parameter at which stability changes sign. Section~5 proves the knee--like crossover in first--passage statistics and quantifies its scaling. Section~6 shows convergence of deterministic exit probabilities to those of the limiting diffusion. Section~7 compares the result with fixed linear Markov descriptions and explains why the knee cannot be reproduced without additional structure. We conclude in Section~8 with implications for projection--based descriptions and abrupt effective transitions more broadly.
\section{A Deterministic Toy Model with Non--Injective Projection}

In this section we define a minimal deterministic dynamical system whose projected dynamics exhibits stochastic behavior and a sharp threshold transition. All ingredients are explicit and finite dimensional. No intrinsic randomness is introduced at the fundamental level.

\subsection{State space and projection}

Let the full state space be the product
\[
\mathcal{M} := \mathbb{R} \times \mathbb{T}^1,
\]
with coordinates $(x,y)$, where $x\in\mathbb{R}$ is the observed (slow) variable and $y\in\mathbb{T}^1:=\mathbb{R}/\mathbb{Z}$ represents hidden (fast) degrees of freedom.

Define the observable projection
\[
P : \mathcal{M} \to \mathbb{R}, \qquad P(x,y)=x.
\]
This projection is non--injective: all points $(x,y)$ with the same $x$ but different $y$ are identified at the observational level.

\subsection{Fast deterministic subsystem}

The hidden variable evolves autonomously under the expanding map
\[
y_{n+1} = T(y_n) := 2y_n \pmod{1}.
\]
The map $T$ is deterministic, invertible modulo null sets, and preserves the Lebesgue probability measure $\mu$ on $\mathbb{T}^1$. It is well known that:
\begin{itemize}
\item $T$ is exponentially mixing for H\"older observables;
\item $T$ satisfies a central limit theorem and a functional invariance principle for zero--mean observables;
\item correlation functions decay exponentially.
\end{itemize}
These properties will be used to derive a diffusive limit for the projected dynamics.

\subsection{Slow subsystem and protocol parameter}

The observed variable $x_n$ evolves according to a weakly coupled fast--slow recurrence:
\begin{equation}
x_{n+1} = x_n + \varepsilon f(x_n;s) + \sqrt{\varepsilon}\,\phi(y_n),
\label{eq:slow-fast-map}
\end{equation}
where:
\begin{itemize}
\item $0<\varepsilon\ll 1$ is a time--scale separation parameter;
\item $s\in\mathbb{R}$ is an external control parameter (``protocol strength'');
\item $f:\mathbb{R}\times\mathbb{R}\to\mathbb{R}$ is a smooth drift function;
\item $\phi:\mathbb{T}^1\to\mathbb{R}$ is a H\"older continuous observable satisfying
\[
\int_{\mathbb{T}^1} \phi(y)\,d\mu(y)=0.
\]
\end{itemize}

A concrete choice used throughout this paper is
\[
\phi(y)=\sin(2\pi y),
\]
which is smooth, bounded, and zero--mean.

\subsection{Linearized drift and stability margin}

We assume that the drift takes the form
\begin{equation}
f(x;s) = -\gamma(s)\,x,
\label{eq:linear-drift}
\end{equation}
where $\gamma:\mathbb{R}\to\mathbb{R}$ is a continuous function of $s$ satisfying:
\begin{enumerate}
\item there exists $s^{\ast}\in\mathbb{R}$ such that $\gamma(s^{\ast})=0$;
\item $\gamma(s)>0$ for $s<s^{\ast}$ (restoring drift);
\item $\gamma(s)<0$ for $s>s^{\ast}$ (repelling drift).
\end{enumerate}

The parameter $\gamma(s)$ will be referred to as the \emph{effective stability margin}. The value $s^{\ast}$ is the critical parameter at which stability of the origin changes sign.

\subsection{Basin boundary and first--passage event}

We define a basin boundary at $x=0$. Initial conditions are chosen with $x_0<0$, corresponding to a metastable basin on the negative half--line.

Define the first--passage (exit) time
\[
\tau := \inf\{n\ge 0 : x_n \ge 0\}.
\]
This event represents a qualitative change in the effective description (``selection'' or ``collapse'').

The primary observable of interest will be the exit probability
\[
p(s,n) := \mathbb{P}_{x_0<0}\bigl(\tau \le n\bigr),
\]
and its continuous--time analogue obtained after rescaling $n=t/\varepsilon$.

\subsection{Remarks}

\begin{enumerate}
\item The full dynamics on $\mathcal{M}$ is deterministic and smooth; no stochastic terms appear at this level.
\item All randomness in the observed variable $x$ arises solely from projection onto $x$ and from the chaotic evolution of the hidden variable $y$.
\item The model is finite dimensional and explicitly solvable up to standard limit theorems, making it suitable for rigorous analysis.
\end{enumerate}

In the next section we derive a diffusive limit for the projected dynamics as $\varepsilon\to 0$ and show that the effective process converges to an Ornstein--Uhlenbeck equation whose stability coefficient is precisely $\gamma(s)$.
\section{Deterministic Homogenization and Diffusive Limit}

In this section we derive a diffusive limit for the projected dynamics defined in
Section~2. The key point is that the effective stochastic behavior of the observed
variable $x$ arises \emph{entirely} from deterministic chaos in the hidden variable
$y$ combined with non--injective projection. No intrinsic randomness is introduced.

\subsection{Rescaled process}

Define the continuous--time interpolation of the discrete process $\{x_n\}$ by
\[
X^\varepsilon(t) := x_{\lfloor t/\varepsilon \rfloor},
\qquad t \ge 0.
\]
We consider the limit $\varepsilon\to 0$ with $t$ fixed.

Substituting \eqref{eq:slow-fast-map} and \eqref{eq:linear-drift} gives
\[
X^\varepsilon(t+\varepsilon) - X^\varepsilon(t)
= \varepsilon f(X^\varepsilon(t);s)
+ \sqrt{\varepsilon}\,\phi(y_{\lfloor t/\varepsilon \rfloor}).
\]

\subsection{Assumptions on the fast subsystem}

The following properties of the fast map $T$ are standard and well established.

\begin{assumption}[Fast mixing and invariance principle]
\label{ass:fast-mixing}
The map $T:\mathbb{T}^1\to\mathbb{T}^1$ defined by $T(y)=2y \pmod{1}$ satisfies:
\begin{enumerate}
\item Exponential decay of correlations for H\"older observables with respect to the invariant measure $\mu$.
\item A functional central limit theorem (invariance principle): for any H\"older observable $\phi$ with zero mean,
\[
\frac{1}{\sqrt{n}}\sum_{k=0}^{n-1} \phi\circ T^k
\;\Rightarrow\;
\sigma W(t),
\]
where $W(t)$ is standard Brownian motion and $\sigma^2>0$ is given by a Green--Kubo formula.
\end{enumerate}
\end{assumption}

These results are classical for expanding maps and can be found, for example, in the literature on deterministic homogenization and dynamical systems.

\subsection{Diffusive limit theorem}

We now state the main homogenization result for the projected dynamics.

\begin{theorem}[Deterministic homogenization for the projected slow variable]
\label{thm:homogenization}
Fix $T>0$ and $s\in\mathbb{R}$. Let $T:\mathbb{T}^1\to\mathbb{T}^1$ be the doubling map
$T(y)=2y\!\!\mod 1$ with invariant Lebesgue measure $\mu$. Let $\phi:\mathbb{T}^1\to\mathbb{R}$
be H\"older and centered, $\int_{\mathbb{T}^1}\phi\,d\mu=0$, and set
\[
S_n := \sum_{k=0}^{n-1}\phi\circ T^k.
\]
Assume the weak invariance principle (functional CLT) holds for $\phi$:
\[
n^{-1/2}S_{\lfloor nt\rfloor}\ \Rightarrow\ \sigma W(t)\quad\text{in }D([0,T])
\]
for some $\sigma>0$ and standard Brownian motion $W$.

Let $f(\cdot;s):\mathbb{R}\to\mathbb{R}$ be globally Lipschitz in $x$, uniformly in $s$ on compact
$s$-sets. Define the deterministic fast--slow recursion
\[
y_{n+1}=T(y_n),\qquad x_{n+1}=x_n+\varepsilon f(x_n;s)+\sqrt{\varepsilon}\,\phi(y_n),
\]
and the rescaled piecewise-constant interpolation $X^\varepsilon(t):=x_{\lfloor t/\varepsilon\rfloor}$
on $[0,T]$.

Then as $\varepsilon\to 0$, $X^\varepsilon \Rightarrow X$ in $D([0,T])$, where $X$ is the unique
weak solution of the SDE
\begin{equation}
dX(t)=f(X(t);s)\,dt+\sigma\,dW(t),\qquad X(0)=x_0.
\label{eq:sde-limit}
\end{equation}
Moreover, the diffusion coefficient $\sigma^2$ is given by the Green--Kubo formula
\begin{equation}
\sigma^2
= \int_{\mathbb{T}^1}\phi(y)^2\,d\mu(y)
+2\sum_{k=1}^\infty\int_{\mathbb{T}^1}\phi(y)\phi(T^k y)\,d\mu(y),
\label{eq:green_kubo_rigorous}
\end{equation}
where the series converges absolutely.
\end{theorem}

\begin{remark}
A fully detailed proof may be obtained by combining the weak invariance principle for the doubling
map with standard fast--slow homogenization arguments for discrete-time recursions; see, e.g.,
Melbourne--Stuart (2011) and Pavliotis--Stuart (2008) for general frameworks.
\end{remark}


\begin{proof}[Sketch of proof]
The result follows from standard deterministic homogenization arguments for fast--slow systems.
Because $y_n$ evolves independently of $x_n$ and satisfies a functional central limit theorem,
the rescaled partial sums of $\phi(y_n)$ converge weakly to Brownian motion.

The slow drift term contributes the deterministic integral
$\int_0^t f(X(s);s)\,ds$ in the limit. Tightness of $X^\varepsilon$ and identification of
the limiting martingale problem yield convergence in distribution to the solution of
\eqref{eq:sde-limit}. A complete proof follows established arguments and is omitted.
\end{proof}

\subsection{Ornstein--Uhlenbeck reduction near the boundary}

Substituting the linearized drift \eqref{eq:linear-drift} into
\eqref{eq:sde-limit} yields, in a neighborhood of the basin boundary $x=0$,
\begin{equation}
dX(t) = -\gamma(s)\,X(t)\,dt + \sigma\, dW(t).
\label{eq:ou}
\end{equation}
Equation \eqref{eq:ou} is the Ornstein--Uhlenbeck process with stability coefficient
$\gamma(s)$.

\begin{remark}
The Ornstein--Uhlenbeck form \eqref{eq:ou} is \emph{derived}, not assumed. It arises as the
universal local approximation of the deterministic projected dynamics near the basin
boundary in the homogenization limit.
\end{remark}

\subsection{Interpretation}

Equation \eqref{eq:ou} makes explicit the role of the stability margin $\gamma(s)$:
\begin{itemize}
\item For $\gamma(s)>0$, the origin is linearly stable and boundary crossing is noise--activated.
\item For $\gamma(s)<0$, the origin is unstable and trajectories are expelled deterministically.
\end{itemize}
The sign change of $\gamma(s)$ at $s=s^{\ast}$ therefore marks a qualitative change in the
effective dynamics. In the next section we show that this change produces a sharp,
finite--strength knee in first--passage statistics.
\section{Finite Critical Point and Stability Regimes}

In this section we identify a finite critical value of the control parameter at which the
stability of the basin boundary changes sign. This prepares the ground for the knee
transition in first--passage statistics.

\subsection{Effective stability margin}

From Section~3, the projected dynamics near the boundary $x=0$ is governed by the
Ornstein--Uhlenbeck equation
\[
dX(t) = -\gamma(s)\,X(t)\,dt + \sigma\, dW(t),
\]
where $\gamma(s)$ is the effective stability margin defined in
\eqref{eq:linear-drift}.

\begin{lemma}[Existence of a finite critical point]
\label{lem:critical-point}
Suppose $\gamma:\mathbb{R}\to\mathbb{R}$ is continuous and satisfies
\[
\gamma(s_1)>0, \qquad \gamma(s_2)<0
\]
for some $s_1<s_2$. Then there exists $s^{\ast}\in(s_1,s_2)$ such that
\[
\gamma(s^{\ast})=0.
\]
\end{lemma}

\begin{proof}
This follows immediately from the intermediate value theorem.
\end{proof}

The value $s^{\ast}$ will be referred to as the \emph{critical protocol strength}.
It is finite and model--dependent, but does not require fine tuning.

\subsection{Qualitative regimes}

The sign of $\gamma(s)$ divides parameter space into two qualitatively distinct regimes:
\begin{itemize}
\item \emph{Stable regime} ($s<s^{\ast}$): $\gamma(s)>0$.  
The origin is linearly stable and trajectories are confined near the basin boundary
unless driven across by fluctuations.
\item \emph{Unstable regime} ($s>s^{\ast}$): $\gamma(s)<0$.  
The origin is linearly unstable and trajectories are deterministically expelled.
\end{itemize}

At $s=s^{\ast}$ the effective linear stability vanishes, and higher--order or stochastic
effects dominate.

In the next section we show that this change of stability produces a sharp knee--like
transition in first--passage statistics.

\section{Knee--Like Transition in First--Passage Statistics}

We now analyze the boundary--crossing statistics of the limiting Ornstein--Uhlenbeck
process and show that the stability change at $s=s^{\ast}$ induces a finite--strength
``knee'' in exit probabilities.

\subsection{First--passage observable}

Let $X(t)$ solve \eqref{eq:ou} with initial condition $X(0)=x_0<0$.
Define the first--passage time
\[
\tau := \inf\{t\ge 0 : X(t)\ge 0\}.
\]
For fixed $t>0$ and parameter $s$, define the exit probability
\[
p(s,t) := \mathbb{P}_{x_0}\bigl(\tau \le t\bigr).
\]

The function $p(s,t)$ is the primary observable used to characterize the knee.

\subsection{Stable regime: noise--activated escape}

Assume $s<s^{\ast}$ so that $\gamma(s)>0$.
In this case the origin is a stable fixed point of the drift.

Standard results for Ornstein--Uhlenbeck processes imply that:
\begin{itemize}
\item boundary crossing occurs only via stochastic fluctuations;
\item for fixed $t$ and small $\sigma$, $p(s,t)$ is exponentially suppressed;
\item mean exit times grow rapidly as $\gamma(s)$ increases.
\end{itemize}

In particular, there exists $C(t,x_0)>0$ such that
\[
p(s,t) \le C(t,x_0)\,
\exp\!\left(-\frac{\gamma(s)x_0^2}{\sigma^2}\right)
\quad \text{for } s<s^{\ast}.
\]

\subsection{Unstable regime: deterministic expulsion}

Assume $s>s^{\ast}$ so that $\gamma(s)<0$.
In this case the drift pushes trajectories toward the boundary.

Elementary estimates for \eqref{eq:ou} yield:
\begin{itemize}
\item exit occurs on a deterministic timescale of order $|\gamma(s)|^{-1}$;
\item noise plays only a subleading role;
\item for fixed $t$ sufficiently large, $p(s,t)\to 1$ as $s$ increases.
\end{itemize}

\subsection{The knee theorem}

We now combine the two regimes into a single crossover statement.

\begin{theorem}[Exact knee for an affine Ornstein--Uhlenbeck readout]
\label{thm:knee}
Assume the limiting SDE in Theorem~\ref{thm:homogenization} has affine drift
\[
f(x;s)=a(s)-\gamma x
\]
with constants $\gamma>0$ and $a:\mathbb{R}\to\mathbb{R}$ continuously differentiable and
strictly increasing in $s$ (so $a'(s)>0$). Let $X(t)$ solve
\begin{equation}
dX(t)=(a(s)-\gamma X(t))\,dt+\sigma\,dW(t),\qquad X(0)=x_0<0,
\label{eq:affine_ou}
\end{equation}
with $\sigma>0$. Fix a readout time $t>0$ and define the selection probability
\[
q(s,t):=\mathbb{P}(X(t)\ge 0).
\]

Then $X(t)$ is Gaussian with mean and variance
\begin{align}
m(s,t) &= x_0e^{-\gamma t}+\frac{a(s)}{\gamma}\bigl(1-e^{-\gamma t}\bigr),\\
v(t) &= \frac{\sigma^2}{2\gamma}\bigl(1-e^{-2\gamma t}\bigr),
\end{align}
and hence
\begin{equation}
q(s,t)=\Phi\!\left(\frac{m(s,t)}{\sqrt{v(t)}}\right),
\label{eq:q_exact}
\end{equation}
where $\Phi$ is the standard normal CDF.

Moreover:
\begin{enumerate}
\item (Unique knee location) There exists a unique $s^{\ast}(t)$ such that $m(s^{\ast}(t),t)=0$, and
$q(s^{\ast}(t),t)=1/2$.
\item (Monotonicity) $q(s,t)$ is strictly increasing in $s$.
\item (Sharpness / knee width) The derivative is
\[
\partial_s q(s,t)=\varphi\!\left(\frac{m(s,t)}{\sqrt{v(t)}}\right)\cdot
\frac{a'(s)}{\gamma}\cdot\frac{1-e^{-\gamma t}}{\sqrt{v(t)}},
\]
where $\varphi$ is the standard normal density. In particular, the maximal slope is attained at
$s=s^{\ast}(t)$ and equals
\[
\max_s \partial_s q(s,t)=
\frac{1}{\sqrt{2\pi}}\cdot
\frac{a'(s^{\ast}(t))}{\gamma}\cdot\frac{1-e^{-\gamma t}}{\sqrt{v(t)}}.
\]
Consequently, a natural knee width scale $\delta s(t)$ (the $s$--interval over which $q$ changes by
an $O(1)$ amount) satisfies
\begin{equation}
\delta s(t)\asymp
\frac{\gamma\sqrt{v(t)}}{a'(s^{\ast}(t))(1-e^{-\gamma t})}
=
\frac{\sigma}{a'(s^{\ast}(t))}\sqrt{\frac{\gamma}{2}}\,
\frac{\sqrt{1-e^{-2\gamma t}}}{1-e^{-\gamma t}}.
\label{eq:knee_width}
\end{equation}
\end{enumerate}
\end{theorem}

\begin{proof}
The explicit solution of \eqref{eq:affine_ou} yields the stated Gaussian law, hence \eqref{eq:q_exact}.
Items (1)--(3) follow by elementary differentiation and the strict monotonicity of $a(s)$.
\end{proof}



\subsection{Interpretation}

Theorem~\ref{thm:knee} shows that the exit probability undergoes a rapid but continuous
change at a finite control parameter.
This ``knee'' is not a phase transition and does not require infinite degrees of freedom.
It is a direct consequence of the loss of linear stability of the effective drift at the
basin boundary.

In the next section we show that this knee behavior is not reproducible by fixed,
protocol--independent linear Markov generators without introducing additional
state--dependent structure.
\section{Convergence of Deterministic Exit Statistics}

In this section we close the logical loop by showing that the exit statistics of the
\emph{original deterministic system} converge to those of the limiting Ornstein--Uhlenbeck
process derived in Section~3. This ensures that the knee phenomenon is not an artifact of
the diffusion approximation, but a genuine feature of the projected deterministic dynamics.

\subsection{Deterministic first--passage time}

Recall the deterministic system defined by
\[
x_{n+1} = x_n + \varepsilon f(x_n;s) + \sqrt{\varepsilon}\,\phi(y_n),
\qquad
y_{n+1} = T(y_n),
\]
with initial condition $x_0<0$.

Define the deterministic first--passage time
\[
\tau^\varepsilon := \inf\{n\ge 0 : x_n \ge 0\}.
\]
We compare $\tau^\varepsilon$ with the first--passage time $\tau$ of the limiting
Ornstein--Uhlenbeck process.

\subsection{Tightness and weak convergence}

From Theorem~\ref{thm:homogenization}, the rescaled process
$X^\varepsilon(t) = x_{\lfloor t/\varepsilon\rfloor}$ converges weakly in
$C([0,T];\mathbb{R})$ to the solution $X(t)$ of the SDE \eqref{eq:ou} for any fixed $T>0$.

Standard tightness arguments for deterministic homogenization imply that the family
$\{X^\varepsilon\}_{\varepsilon>0}$ is tight and that the convergence holds jointly with
the corresponding driving processes.

\subsection{Convergence of exit events}

We now relate weak convergence of paths to convergence of exit probabilities.

\begin{lemma}[Continuity of the exit functional]
\label{lem:exit-continuity}
Let $\Gamma:C([0,T];\mathbb{R})\to\{0,1\}$ be the exit functional
\[
\Gamma(\omega)=\mathbf{1}\{\exists\, t\in[0,T]: \omega(t)\ge 0\}.
\]
Then $\Gamma$ is continuous at any path that does not touch the boundary $0$ tangentially.
\end{lemma}

For the Ornstein--Uhlenbeck process \eqref{eq:ou}, the probability of tangential boundary
touching is zero. Consequently, $\Gamma$ is almost surely continuous with respect to the
law of $X(t)$.

\begin{theorem}[Convergence of deterministic readout probabilities]
\label{thm:readout_convergence}
Assume Theorem~\ref{thm:homogenization}. Fix $t>0$ such that $X$ has a continuous
distribution at time $t$ (in particular, for the SDE \eqref{eq:affine_ou} with $\sigma>0$, this holds
for all $t>0$). Define the deterministic readout probability
\[
q^\varepsilon(s,t):=\mathbb{P}\bigl(X^\varepsilon(t)\ge 0\bigr)
\qquad\text{and}\qquad
q(s,t):=\mathbb{P}\bigl(X(t)\ge 0\bigr).
\]
Then
\[
\lim_{\varepsilon\to 0} q^\varepsilon(s,t)=q(s,t).
\]
\end{theorem}

\begin{proof}
By Theorem~\ref{thm:homogenization}, $X^\varepsilon \Rightarrow X$ in $D([0,T])$.
The evaluation map $\pi_t:D([0,T])\to\mathbb{R}$, $\pi_t(\omega)=\omega(t)$, is continuous at any
path that is continuous at time $t$. Since $X$ has continuous sample paths almost surely, $\pi_t$ is
$X$-a.s.\ continuous. Hence by the continuous mapping theorem, $X^\varepsilon(t)\Rightarrow X(t)$.

Now consider the indicator $g(z)=\mathbf{1}\{z\ge 0\}$. The function $g$ is continuous except at
$z=0$. Since $X(t)$ has a continuous distribution, $\mathbb{P}(X(t)=0)=0$, and by the
Portmanteau theorem (or a standard lemma on convergence of expectations for bounded functions
continuous $\nu$-a.s.), we obtain
\[
\lim_{\varepsilon\to 0}\mathbb{E}\,g(X^\varepsilon(t))=\mathbb{E}\,g(X(t)),
\]
i.e.\ $q^\varepsilon(s,t)\to q(s,t)$.
\end{proof}



\subsection{Persistence of the knee}

Combining Theorem~\ref{thm:knee} with Theorem~\ref{thm:readout_convergence}, we obtain:

\begin{corollary}[Deterministic knee]
\label{cor:deterministic-knee}
For fixed $t>0$, the deterministic exit probability
\[
p^\varepsilon(s,t)
:=
\mathbb{P}\!\left(\tau^\varepsilon \le \frac{t}{\varepsilon}\right)
\]
exhibits a sharp crossover near $s=s^{\ast}$ for sufficiently small $\varepsilon$.
The location and scaling of the knee converge to those of the limiting
Ornstein--Uhlenbeck process as $\varepsilon\to 0$.
\end{corollary}

\subsection{Interpretation}

Corollary~\ref{cor:deterministic-knee} shows that the knee--like transition is already
present at the level of the deterministic dynamics once the system is viewed through the
non--injective projection $P(x,y)=x$. The diffusion approximation does not create the knee;
it merely renders the mechanism analytically transparent.

This establishes that sharp effective threshold behavior can arise from smooth,
invertible dynamics without introducing intrinsic stochasticity or discontinuous rules.
\section{Comparison with Fixed Linear Markov Descriptions}

In this section we show that the knee--like transition exhibited by the projected
deterministic dynamics cannot be reproduced by fixed, protocol--independent linear
Markov models without introducing additional state--dependent structure. This
distinguishes the present mechanism from standard decoherence--only or
phenomenological collapse descriptions.

\subsection{Fixed linear Markov models}

Consider an effective description of the observed variable $x$ given by a fixed
linear Markov semigroup with generator $L$, independent of the control parameter $s$.
In continuous time this takes the form
\[
\partial_t \rho(t) = L \rho(t),
\]
where $\rho(t)$ is a probability density on $\mathbb{R}$ and $L$ is a second--order
differential operator with smooth coefficients.

Typical examples include:
\begin{itemize}
\item Ornstein--Uhlenbeck generators with fixed drift and diffusion;
\item Fokker--Planck equations with constant coefficients;
\item Lindblad--type generators in continuous measurement models with fixed rates.
\end{itemize}

In all such cases, the qualitative behavior of exit statistics is governed by the
spectrum and coefficients of $L$.

\subsection{Absence of finite--strength knees}

For a fixed generator $L$, exit rates and first--passage statistics depend smoothly on
initial conditions and observation time. In particular:
\begin{itemize}
\item there is no finite control parameter at which the sign of the drift changes unless
      $L$ itself is modified;
\item any sharp change in exit behavior must be encoded explicitly in the generator;
\item protocol dependence can only enter through time--dependent or state--dependent
      modification of $L$.
\end{itemize}

Consequently, a fixed linear Markov model cannot exhibit a finite--strength knee
transition of the form established in Theorem~\ref{thm:knee} unless the generator is
allowed to depend on the control parameter or on hidden state variables.

\subsection{Structural origin of the knee}

By contrast, in the deterministic model of Sections~2--6:
\begin{itemize}
\item the ``generator'' of the effective dynamics is not fixed a priori;
\item the effective drift coefficient $\gamma(s)$ emerges from the interaction between
      the observed and hidden variables;
\item the sign change of $\gamma(s)$ is forced by loss of stability in the projected
      description, not by ad hoc modification.
\end{itemize}

The knee therefore reflects a structural instability of the effective encoding induced
by non--injective projection, rather than a feature inserted at the level of the reduced
dynamics.

\subsection{Implications}

This comparison shows that the knee--like transition derived in this paper is not a
generic feature of Markovian noise models, but a consequence of state--dependent
effective dynamics generated by projection. Any framework that treats the effective
generator as fixed must either fail to reproduce the knee or introduce additional
mechanisms by hand.

\section{Discussion and Conclusion}

\subsection{Summary of results}

We have constructed a fully explicit, finite--dimensional deterministic dynamical system
in which:
\begin{enumerate}
\item the full dynamics is smooth, deterministic, and invertible;
\item the observed variable is obtained by a non--injective projection;
\item deterministic homogenization yields a rigorous diffusive limit;
\item a finite control parameter produces a sharp knee--like crossover in exit statistics.
\end{enumerate}

The knee arises when an effective stability margin crosses zero, separating a
noise--activated escape regime from deterministic expulsion. The phenomenon is
continuous but sharply bent, requires no thermodynamic limit, and persists in the
original deterministic system.

\subsection{Interpretational significance}

The result demonstrates that abrupt effective transitions---often associated with
intrinsic randomness or discontinuous rules---can arise purely from projection and
stability loss. The knee is not an artifact of stochastic modeling but a structural
consequence of reduced description.

This provides a concrete benchmark against which projection--based and
decoherence--only descriptions can be compared. In particular, it shows that fixed
linear Markov generators are insufficient to capture protocol--dependent threshold
behavior without additional structure.

\subsection{Outlook}

While the present work is intentionally minimal, the mechanism extends naturally to
higher--dimensional systems, open quantum models, and more elaborate projection
schemes. The explicit knee derived here can serve as a stress test for any framework
that claims to explain collapse--like or threshold phenomena without introducing
ad hoc dynamics.

\subsection*{Acknowledgements}

The author thanks colleagues for discussions on deterministic homogenization,
metastability, and projection--based descriptions.
\bigskip
\nocite{*}
\bibliographystyle{plain}
\bibliography{main}

\end{document}
